% =============================================================================
% Kapitel 5: Fazit und Handlungsempfehlungen
% =============================================================================

\section{Fazit und Handlungsempfehlungen}
\label{sec:fazit}

Die Fallstudie hat fünf Failure Cases des ROSI-Systems auf Prozess- und
Betriebssystemebene untersucht und damit die drei eingangs formulierten
Forschungsfragen beantwortet.

Zu FF2 (Systembeobachtbarkeit) zeigt Case~1 die gravierendste Lücke: Drei von vier
Service-Prozessen fallen ohne Erkennung aus. Da das System über kein externes Alerting
verfügt, bleibt die Time-to-Detection in diesen Fällen unbegrenzt, was sich unmittelbar
negativ auf die Wiederherstellungszeit und damit die Verfügbarkeit nach
Gleichung~\ref{eq:availability} auswirkt. Zu FF1 (Ressourcenverwaltung) belegt Case~2 die
fehlende Rückfallebene des flüchtigen Bildspeichers, während Case~5 verdeutlicht, dass
nicht jede plausible Schwachstelle auch real ist. Zu FF3 (Datenpersistenz und Kaskaden)
zeigen Case~3 und Case~4 zwei unterschiedliche Mechanismen: einen systemweiten Freeze durch
Backpressure sowie ein durch die Entwicklungskonfiguration maskiertes,
produktionsspezifisches Verbindungsrisiko.

\subsection{Übergreifende Befunde}
\label{sec:fazit-uebergreifend}

Zwei Muster ziehen sich durch die Ergebnisse. Erstens fehlt durchgängig eine
\textit{Graceful Degradation}: Ein lokaler Fehler -- ein abgestürzter Prozess, ein
fehlgeschlagener Schreibvorgang, eine tote Verbindung -- führt regelmäßig zum Ausfall des
Gesamtsystems statt zu einem kontrollierten Teilbetrieb. Zweitens bestehen
\textit{blinde Flecken im Monitoring}: Weder Queue-Tiefen noch Speicherfüllstände noch der
Zustand der Service-Prozesse werden zur Laufzeit überwacht. Die Cases~3, 4 und 5 stehen
zudem in einem Kaskadenzusammenhang, da eine verlangsamte Datenbank über die Log-Queue das
gesamte System ausbremsen und so weitere Fehlerzustände begünstigen kann.

\subsection{Priorisierte Handlungsempfehlungen}
\label{sec:fazit-empfehlungen}

Tabelle~\ref{tab:empfehlungen} fasst die abgeleiteten Maßnahmen zusammen. Die ersten vier
Maßnahmen greifen direkt in die Verfügbarkeit ein, weil sie entweder die
Time-to-Detection verkürzen oder einen lokalen Fehler vom Gesamtsystem entkoppeln. Case~5
erfordert dagegen keinen Sofortfix, sollte aber als Monitoring-Indikator für langsame
Bildspeicherung genutzt werden.

\begin{table}[htbp]
    \centering
    \caption{Priorisierte Handlungsempfehlungen}
    \label{tab:empfehlungen}
    \begin{tabularx}{\textwidth}{llX}
        \hline
        \textbf{Case} & \textbf{Aufwand} & \textbf{Maßnahme} \\
        \hline
        1 & gering   & Supervision aller Service-Prozesse oder \texttt{systemd}-Unit \\
        2 & gering   & Fehlerbehandlung in der Bildspeicherung + Füllstands-Monitoring \\
        3 & minimal  & \texttt{put(timeout)} mit Drop-on-Full in der Log-Queue \\
        4 & minimal  & \texttt{close\_old\_connections()} bzw. \texttt{CONN\_MAX\_AGE = 0} \\
        5 & --       & kein Sofortfix; optional Referenzzähler und Step-5-Monitoring \\
        \hline
    \end{tabularx}
\end{table}

\subsection{Ausblick}
\label{sec:fazit-ausblick}

Die größten Hebel für die Produktionsreife liegen in einer durchgängigen
Laufzeit-Supervision aller Prozesse sowie einem Monitoring der bislang unbeobachteten
Ressourcen. Beide Maßnahmen adressieren nicht nur einzelne Cases, sondern die
übergreifenden Muster fehlender Beobachtbarkeit und fehlender Graceful Degradation. Eine
Übertragung der Untersuchung auf das spätere Kundensystem -- mit abweichender Hardware und
aktiver Produktionskonfiguration -- bleibt offen und würde insbesondere das in Case~4
analytisch hergeleitete Verbindungsrisiko empirisch absichern.
